\section{Open rank cells of finite open chains}

\import{topology/urysohn/chain.tex}
\import{topology/urysohn/rank.tex}
\import{logic/conditional.tex}
\import{set/logical-graph.tex}

To control step values locally, we cover the space by open sets on which
entry ranks differ by at most one. An interior cell lies in the next chain
entry and outside the closure of the preceding entry. Closure nesting makes
these cells open and ensures that they cover the space together with the two
endpoint cells.

The formulas specify the endpoints separately, so no predecessor is needed
at zero. They define total terms; the open-cover and rank properties below
require a chain with a positive number of steps.

\begin{definition}\label{open_chain_rank_cell}
    $\openChainRankCell{X}{c}{r}{k}=
        \Apply{\logicalIfSetOp}{(k=\emptyset),c(\suc{\emptyset}),
            \Apply{\logicalIfSetOp}{(k=r),\carrier[X]\setminus
                    \closure{c(\unions{r})}{X},
                c(\suc{k})\setminus
                    \closure{c(\unions{k})}{X}}}$.
\end{definition}

\begin{definition}\label{open_chain_rank_cells}
    $\openChainRankCells{X}{c}{r}=
        \Apply{\graphOnDomainOp}
            {\suc{r},
                \fun{k:\Set}{\openChainRankCell{X}{c}{r}{k}}}$.
\end{definition}

\begin{proposition}\label{open_chain_rank_cell_zero}
    $\openChainRankCell{X}{c}{r}{\emptyset}=c(\suc{\emptyset})$.
\end{proposition}
\begin{proof}
    Follows by
        \obv{open_chain_rank_cell,logical_if_set_true}.
\end{proof}

\begin{proposition}\label{open_chain_rank_cell_last}
    Suppose $r\neq\emptyset$.
    Then
        $\openChainRankCell{X}{c}{r}{r}
            =\carrier[X]\setminus\closure{c(\unions{r})}{X}$.
\end{proposition}
\begin{proof}
    Follows by
        \auto{open_chain_rank_cell,logical_if_set_false,logical_if_set_true}.
\end{proof}

\begin{proposition}\label{open_chain_rank_cell_interior}
    Suppose $k\neq\emptyset$ and $k\neq r$.
    Then
        $\openChainRankCell{X}{c}{r}{k}
            =c(\suc{k})\setminus\closure{c(\unions{k})}{X}$.
\end{proposition}
\begin{proof}
    Follows by
        \auto{open_chain_rank_cell,logical_if_set_false}.
\end{proof}

\begin{corollary}\label{open_chain_rank_cells_range_finite}
    Let $r$ be a natural number.
    Then $\ran{\openChainRankCells{X}{c}{r}}$ is finite.
\end{corollary}
\begin{proof}
    Let $D=\openChainRankCells{X}{c}{r}$.
    $D\in\funs{\suc{r}}{\ran{D}}$
        by \auto{open_chain_rank_cells,logical_graph_on_domain_function}.
    Hence $D$ is a function and $\dom{D}=\suc{r}$
        by \obv{funs_is_function,funs}.
    $\suc{r}$ is a natural number by \obv{omega_successor}.
    Hence $\suc{r}$ is finite by \obv{natural_number_is_finite}.
    $\suc{r}\subseteq\dom{D}$ by \obv{subseteq_refl}.
    Hence $\img{D}{\suc{r}}$ is finite
        by \obv{finite_function_image}.
    $\img{D}{\suc{r}}=\ran{D}$ by \auto{img_dom}.
\end{proof}

\begin{proposition}\label{open_chain_rank_cells_apply}
    Let $k\in\suc{r}$.
    Then $\apply{\openChainRankCells{X}{c}{r}}{k}
        =\openChainRankCell{X}{c}{r}{k}$.
\end{proposition}
\begin{proof}
    Follows by
        \auto{open_chain_rank_cells,logical_graph_on_domain_apply}.
\end{proof}

\begin{corollary}\label{open_chain_rank_cells_zero}
    Suppose $\emptyset\in\suc{r}$.
    Then $\apply{\openChainRankCells{X}{c}{r}}{\emptyset}
        =c(\suc{\emptyset})$.
\end{corollary}
\begin{proof}
    Follows by
        \auto{open_chain_rank_cells_apply,open_chain_rank_cell_zero}.
\end{proof}

\begin{corollary}\label{open_chain_rank_cells_last}
    Suppose $r\neq\emptyset$.
    Then $\apply{\openChainRankCells{X}{c}{r}}{r}
        =\carrier[X]\setminus\closure{c(\unions{r})}{X}$.
\end{corollary}
\begin{proof}
    $r\in\suc{r}$ by \obv{suc_intro_self}.
    Follows by
        \auto{open_chain_rank_cells_apply,open_chain_rank_cell_last}.
\end{proof}

\begin{corollary}\label{open_chain_rank_cells_interior}
    Let $k\in r$.
    Suppose $k\neq\emptyset$.
    Then $\apply{\openChainRankCells{X}{c}{r}}{k}
        =c(\suc{k})\setminus\closure{c(\unions{k})}{X}$.
\end{corollary}
\begin{proof}
    $k\in\suc{r}$ by \obv{suc_intro_in}.
    $k\neq r$ by \obv{in_irrefl}.
    Follows by
        \auto{open_chain_rank_cells_apply,open_chain_rank_cell_interior}.
\end{proof}

\begin{proposition}\label{open_chain_rank_cell_open}
    Let $X$ be a topological space.
    Let $r$ be a natural number such that $r\neq\emptyset$.
    Suppose $c$ is a finite open chain with $r$ steps in $X$.
    Let $k\in\suc{r}$.
    Then
        $\openChainRankCell{X}{c}{r}{k}
            \in\opens[X]$.
\end{proposition}
\begin{proof}
    $c\in\funs{\suc{r}}{\opens[X]}$
        by \obv{finite_open_chain_type}.
    \begin{byCase}
        \caseOf{$k=\emptyset$.}
            $\emptyset\in r$ by \obv{natural_nonzero_contains_zero}.
            $\emptyset$ is a natural number by \obv{omega_empty}.
            $\suc{\emptyset}\in\suc{r}$
                by \obv{natural_successor_in_successor_bound}.
            Hence $c(\suc{\emptyset})\in\opens[X]$
                by \obv{funs_type_apply}.
            Follows by \auto{open_chain_rank_cell_zero}.
        \caseOf{$k\neq\emptyset$.}
            \begin{byCase}
                \caseOf{$k=r$.}
                    $\closure{c(\unions{r})}{X}$ is closed in $X$
                        by \obv{closure_is_closed}.
                    Hence $\closure{c(\unions{r})}{X}
                        \subseteq\carrier[X]$
                        by \obv{closed_in_subseteq_carrier}.
                    Thus $\carrier[X]\setminus
                        \closure{c(\unions{r})}{X}$ is open in $X$
                        by \obv{complement_of_closed_is_open}.
                    Hence $\carrier[X]\setminus
                        \closure{c(\unions{r})}{X}\in\opens[X]$
                        by \obv{is_open}.
                    Follows by \auto{open_chain_rank_cell_last}.
                \caseOf{$k\neq r$.}
                    Hence $k\in r$ by \auto{suc_iff}.
                    $r\subseteq\naturals$
                        by
                        \obv{naturals_is_transitive_set,transitiveset_iff_subseteq}.
                    $\suc{k}\in\suc{r}$
                        by \obv{natural_successor_in_successor_bound}.
                    Hence $c(\suc{k})\in\opens[X]$
                        by \obv{funs_type_apply}.
                    Thus $c(\suc{k})$ is open in $X$ by \obv{is_open}.
                    Hence $c(\suc{k})\subseteq\carrier[X]$
                        by \obv{opens_type,subseteq}.
                    $\closure{c(\unions{k})}{X}$ is closed in $X$
                        by \obv{closure_is_closed}.
                    Hence $\closure{c(\unions{k})}{X}
                        \subseteq\carrier[X]$
                        by \obv{closed_in_subseteq_carrier}.
                    Thus $c(\suc{k})\setminus
                        \closure{c(\unions{k})}{X}$ is open in $X$
                        by \obv{opens_minus_closed_is_open}.
                    Hence $c(\suc{k})\setminus
                        \closure{c(\unions{k})}{X}\in\opens[X]$
                        by \obv{is_open}.
                    Follows by
                        \obv{open_chain_rank_cell_interior}.
            \end{byCase}
    \end{byCase}
\end{proof}

\begin{theorem}\label{open_chain_rank_cells_typed}
    Let $X$ be a topological space.
    Let $r$ be a natural number such that $r\neq\emptyset$.
    Suppose $c$ is a finite open chain with $r$ steps in $X$.
    Then $\openChainRankCells{X}{c}{r}
        \in\funs{\suc{r}}{\opens[X]}$.
\end{theorem}
\begin{proof}
    For every $k\in\suc{r}$ we have
        $\openChainRankCell{X}{c}{r}{k}\in\opens[X]$
        by \obv{open_chain_rank_cell_open}.
    Follows by
        \auto{open_chain_rank_cells,logical_graph_on_domain_type}.
\end{proof}

\subsection{Ranks on the cells}

\begin{lemma}\label{open_chain_entry_rank_after_closed_entry}
    Let $X$ be a topological space.
    Suppose $c$ is a finite open chain with $r$ steps in $X$.
    Let $p\in r$.
    Suppose $x\notin\closure{c(p)}{X}$.
    Then $p\in\entryRank{c}{r}{x}$.
\end{lemma}
\begin{proof}
    $r$ is a natural number and
        $c\in\funs{\suc{r}}{\opens[X]}$
        by \obv{finite_open_chain_type}.
    $p$ is a natural number
        by
        \obv{naturals_is_transitive_set,transitiveset_iff_subseteq,subseteq}.
    Let $q=\entryRank{c}{r}{x}$.
    $q$ is an ordinal by \obv{entry_rank_is_ordinal}.
    $p$ is an ordinal by \obv{natural_number_is_ordinal}.
    Suppose not.
    $q\in\entryRankCandidates{c}{r}{x}$
        by \auto{entry_rank_is_candidate}.
    Hence $q\in r$ and $x\in c(q)$ or $q=r$
        by \obv{entry_rank_candidates_iff}.
    Show that $q\in r$ and $x\in c(q)$.
    \begin{subproof}
        Suppose not.
        Hence $q=r$.
        Then $p\in q$.
        Contradiction.
    \end{subproof}
    $q\in p$ or $q=p$ or $p\in q$
        by \obv{ordinal_elem_connex}.
    Hence $q\in p$ or $q=p$.
    Hence $q\subseteq p$
        by
        \obv{ordinal_elem_implies_subseteq,subseteq_refl,natural_number_is_ordinal}.
    $q\in\suc{r}$ and $p\in\suc{r}$ by \obv{suc_intro_in}.
    Thus $c(q)\subseteq c(p)$
        by \obv{finite_open_chain_monotone}.
    $c(p)\in\opens[X]$ by \obv{funs_type_apply,suc_intro_in}.
    Hence $c(p)\subseteq\carrier[X]$
        by \obv{opens_type,subseteq}.
    Thus $c(p)\subseteq\closure{c(p)}{X}$
        by \obv{set_is_subseteq_to_closure_of_the_set}.
    Hence $c(q)\subseteq\closure{c(p)}{X}$
        by \obv{subseteq_transitive}.
    Thus $x\in\closure{c(p)}{X}$ by \obv{subseteq}.
    Contradiction.
\end{proof}

\begin{proposition}\label{open_chain_entry_rank_on_zero_cell}
    Let $X$ be a topological space.
    Let $r$ be a natural number such that $r\neq\emptyset$.
    Suppose $c$ is a finite open chain with $r$ steps in $X$.
    Let $x\in\apply{\openChainRankCells{X}{c}{r}}{\emptyset}$.
    Then $\entryRank{c}{r}{x}=\emptyset$ or
        $\entryRank{c}{r}{x}=\suc{\emptyset}$.
\end{proposition}
\begin{proof}
    $\emptyset\in r$ by \obv{natural_nonzero_contains_zero}.
    Hence $\emptyset\in\suc{r}$ by \obv{suc_intro_in}.
    $x\in c(\suc{\emptyset})$
        by \obv{open_chain_rank_cells_zero}.
    $\emptyset$ is a natural number by \obv{omega_empty}.
    Hence $\suc{\emptyset}\in\suc{r}$
        by \obv{natural_successor_in_successor_bound}.
    Thus $\suc{\emptyset}\in r$ or
        $\suc{\emptyset}=r$ by \auto{suc_iff}.
    Hence $\suc{\emptyset}
        \in\entryRankCandidates{c}{r}{x}$
        by \auto{entry_rank_candidates_iff}.
    Thus $\entryRank{c}{r}{x}\subseteq\suc{\emptyset}$
        by \obv{entry_rank_subseteq}.
    $\entryRank{c}{r}{x}$ and $\suc{\emptyset}$ are ordinals
        by
        \obv{entry_rank_is_ordinal,emptyset_is_ordinal,suc_ordinal}.
    Hence $\entryRank{c}{r}{x}\in\suc{\emptyset}$ or
        $\entryRank{c}{r}{x}=\suc{\emptyset}$
        by \obv{ordinal_subseteq_implies_elem_or_eq}.
    Follows by \auto{suc_iff,emptyset}.
\end{proof}

\begin{proposition}\label{open_chain_entry_rank_on_interior_cell}
    Let $X$ be a topological space.
    Let $r$ be a natural number such that $r\neq\emptyset$.
    Suppose $c$ is a finite open chain with $r$ steps in $X$.
    Let $k\in r$.
    Suppose $k\neq\emptyset$.
    Let $x\in\apply{\openChainRankCells{X}{c}{r}}{k}$.
    Then $\entryRank{c}{r}{x}=k$ or
        $\entryRank{c}{r}{x}=\suc{k}$.
\end{proposition}
\begin{proof}
    Let $p=\unions{k}$.
    $r\subseteq\naturals$
        by \obv{naturals_is_transitive_set,transitiveset_iff_subseteq}.
    Hence $k$ is a natural number by \obv{subseteq}.
    $p$ is a natural number and $k=\suc{p}$
        by \auto{natural_nonzero_successor_of_union}.
    $p\in k$ by \obv{suc_intro_self}.
    Hence $p\in r$
        by
        \obv{natural_number_is_ordinal,ordinal_transitivity}.
    $x\in c(\suc{k})$ and
        $x\notin\closure{c(p)}{X}$
        by
        \auto{open_chain_rank_cells_interior,setminus}.
    Hence $p\in\entryRank{c}{r}{x}$
        by \obv{open_chain_entry_rank_after_closed_entry}.
    $\entryRank{c}{r}{x}$ is an ordinal
        by \obv{entry_rank_is_ordinal}.
    Hence $\suc{p}\subseteq\entryRank{c}{r}{x}$
        by \obv{ordinal_suc_subseteq}.
    Thus $k\subseteq\entryRank{c}{r}{x}$.
    $\suc{k}$ is a natural number by \obv{omega_successor}.
    Hence $\suc{k}$ is an ordinal by \obv{natural_number_is_ordinal}.
    $\suc{k}\in\suc{r}$
        by \obv{natural_successor_in_successor_bound}.
    Thus $\suc{k}\in r$ or $\suc{k}=r$ by \auto{suc_iff}.
    Hence $\suc{k}\in\entryRankCandidates{c}{r}{x}$
        by \auto{entry_rank_candidates_iff}.
    Thus $\entryRank{c}{r}{x}\subseteq\suc{k}$
        by \obv{entry_rank_subseteq}.
    Hence $\entryRank{c}{r}{x}\in\suc{k}$ or
        $\entryRank{c}{r}{x}=\suc{k}$
        by \obv{ordinal_subseteq_implies_elem_or_eq}.
    \begin{byCase}
        \caseOf{$\entryRank{c}{r}{x}=\suc{k}$.}
            Follows.
        \caseOf{$\entryRank{c}{r}{x}\neq\suc{k}$.}
            Hence $\entryRank{c}{r}{x}\in\suc{k}$.
            Thus $\entryRank{c}{r}{x}\in k$ or
                $\entryRank{c}{r}{x}=k$ by \auto{suc_iff}.
            Show that $\entryRank{c}{r}{x}=k$.
            \begin{subproof}
                Suppose not.
                Hence $\entryRank{c}{r}{x}\in k$.
                Hence $\entryRank{c}{r}{x}
                    \in\entryRank{c}{r}{x}$
                    by \obv{subseteq}.
                Contradiction by \obv{in_irrefl}.
            \end{subproof}
    \end{byCase}
\end{proof}

\begin{proposition}\label{open_chain_entry_rank_on_last_cell}
    Let $X$ be a topological space.
    Let $r$ be a natural number such that $r\neq\emptyset$.
    Suppose $c$ is a finite open chain with $r$ steps in $X$.
    Let $x\in\apply{\openChainRankCells{X}{c}{r}}{r}$.
    Then $\entryRank{c}{r}{x}=r$.
\end{proposition}
\begin{proof}
    Let $p=\unions{r}$.
    $p$ is a natural number and $r=\suc{p}$
        by \auto{natural_nonzero_successor_of_union}.
    $p\in r$ by \obv{suc_intro_self}.
    $x\notin\closure{c(p)}{X}$
        by \auto{open_chain_rank_cells_last,setminus}.
    Hence $p\in\entryRank{c}{r}{x}$
        by \obv{open_chain_entry_rank_after_closed_entry}.
    $\entryRank{c}{r}{x}$ is an ordinal
        by \obv{entry_rank_is_ordinal}.
    Hence $\suc{p}\subseteq\entryRank{c}{r}{x}$
        by \obv{ordinal_suc_subseteq}.
    Thus $r\subseteq\entryRank{c}{r}{x}$.
    $\entryRank{c}{r}{x}\subseteq r$
        by \obv{entry_rank_subseteq_sentinel}.
    Follows by \obv{subseteq_antisymmetric}.
\end{proof}

\subsection{The open cover}

\begin{theorem}\label{open_chain_rank_cells_cover_point}
    Let $X$ be a topological space.
    Let $r$ be a natural number such that $r\neq\emptyset$.
    Suppose $c$ is a finite open chain with $r$ steps in $X$.
    Let $x\in\carrier[X]$.
    Then there exists $k\in\suc{r}$ such that
        $x\in\apply{\openChainRankCells{X}{c}{r}}{k}$.
\end{theorem}
\begin{proof}
    $c\in\funs{\suc{r}}{\opens[X]}$
        by \obv{finite_open_chain_type}.
    \begin{byCase}
        \caseOf{$x\in c(r)$.}
            Let $q=\entryRank{c}{r}{x}$.
            $q$ is a natural number by \auto{entry_rank_is_natural}.
            $q=\emptyset$ or there exists a natural number $j$ such that
                $q=\suc{j}$ by \obv{natural_zero_or_successor}.
            \begin{byCase}
                \caseOf{$q=\emptyset$.}
                    $\emptyset\in r$ and $x\in c(\emptyset)$
                        by \auto{entry_rank_eq_zero_iff}.
                    Hence $c(\emptyset)\subseteq c(\suc{\emptyset})$
                        by \obv{finite_open_chain_adjacent}.
                    Thus $x\in c(\suc{\emptyset})$ by \obv{subseteq}.
                    $\emptyset\in\suc{r}$ by \obv{suc_intro_in}.
                    Hence
                        $x\in
                            \apply{\openChainRankCells{X}{c}{r}}{\emptyset}$
                        by \obv{open_chain_rank_cells_zero}.
                \caseOf{$q\neq\emptyset$.}
                    Take a natural number $j$ such that $q=\suc{j}$.
                    $j\in r$
                        by \obv{entry_rank_successor_predecessor}.
                    $\suc{j}=r$ or
                        $\suc{j}\in r$ and $x\in c(\suc{j})$
                        by \obv{entry_rank_successor_candidate}.
                    Hence $x\in c(\suc{j})$.
                    \begin{byCase}
                        \caseOf{$j=\emptyset$.}
                            $j\in\suc{r}$ by \obv{suc_intro_in}.
                            Hence
                                $x\in\apply{
                                    \openChainRankCells{X}{c}{r}}{j}$
                                by
                                \auto{open_chain_rank_cells_zero}.
                        \caseOf{$j\neq\emptyset$.}
                            Let $p=\unions{j}$.
                            $p$ is a natural number and $j=\suc{p}$
                                by
                                \auto{natural_nonzero_successor_of_union}.
                            $p\in j$ by \obv{suc_intro_self}.
                            Hence $p\in r$
                                by
                                \obv{natural_number_is_ordinal,ordinal_transitivity}.
                            $\closure{c(p)}{X}\subseteq c(\suc{p})$
                                by \obv{finite_open_chain_nesting}.
                            Hence $\closure{c(p)}{X}\subseteq c(j)$.
                            $j\in q$ by \obv{suc_intro_self}.
                            Thus $x\notin c(j)$
                                by \obv{entry_rank_no_earlier_entry}.
                            Hence $x\notin\closure{c(p)}{X}$
                                by \obv{subseteq}.
                            Thus
                                $x\in c(\suc{j})\setminus
                                    \closure{c(p)}{X}$
                                by \obv{setminus}.
                            Hence
                                $x\in\apply{
                                    \openChainRankCells{X}{c}{r}}{j}$
                                by
                                \auto{open_chain_rank_cells_interior}.
                            $j\in\suc{r}$ by \obv{suc_intro_in}.
                    \end{byCase}
            \end{byCase}
        \caseOf{$x\notin c(r)$.}
            Let $p=\unions{r}$.
            $p$ is a natural number and $r=\suc{p}$
                by \auto{natural_nonzero_successor_of_union}.
            $p\in r$ by \obv{suc_intro_self}.
            $\closure{c(p)}{X}\subseteq c(\suc{p})$
                by \obv{finite_open_chain_nesting}.
            Hence $\closure{c(p)}{X}\subseteq c(r)$.
            Thus $x\notin\closure{c(p)}{X}$ by \obv{subseteq}.
            Hence
                $x\in\carrier[X]\setminus\closure{c(p)}{X}$
                by \obv{setminus}.
            Thus $x\in\apply{\openChainRankCells{X}{c}{r}}{r}$
                by \auto{open_chain_rank_cells_last}.
            $r\in\suc{r}$ by \obv{suc_intro_self}.
    \end{byCase}
\end{proof}

\begin{corollary}\label{open_chain_rank_cells_cover}
    Let $X$ be a topological space.
    Let $r$ be a natural number such that $r\neq\emptyset$.
    Suppose $c$ is a finite open chain with $r$ steps in $X$.
    Then
        $\unions{\ran{\openChainRankCells{X}{c}{r}}}
            =\carrier[X]$.
\end{corollary}
\begin{proof}
    Let $D=\openChainRankCells{X}{c}{r}$.
    $D\in\funs{\suc{r}}{\opens[X]}$
        by \auto{open_chain_rank_cells_typed}.
    $\ran{D}\subseteq\opens[X]$ by \obv{funs_ran}.
    Show that for every $x\in\unions{\ran{D}}$ we have
        $x\in\carrier[X]$.
    \begin{subproof}
        Fix $x\in\unions{\ran{D}}$.
        Take $V\in\ran{D}$ such that $x\in V$
            by \obv{unions_iff}.
        $V\in\opens[X]$ by \obv{subseteq}.
        Hence $V\subseteq\carrier[X]$
            by \obv{opens_type,subseteq}.
        Follows by \obv{subseteq}.
    \end{subproof}
    Hence $\unions{\ran{D}}\subseteq\carrier[X]$
        by \obv{subseteq}.
    Show that for every $x\in\carrier[X]$ we have
        $x\in\unions{\ran{D}}$.
    \begin{subproof}
        Fix $x\in\carrier[X]$.
        Take $k\in\suc{r}$ such that $x\in D(k)$
            by \auto{open_chain_rank_cells_cover_point}.
        $\suc{r}=\dom{D}$ by \obv{funs}.
        Hence $k\in\dom{D}$.
        $D$ is a function by \obv{funs_is_function}.
        Hence $(k,D(k))\in D$ by \obv{function_apply_elim}.
        Thus $D(k)\in\ran{D}$ by \auto{ran_iff}.
        Follows by \auto{unions_iff}.
    \end{subproof}
    Hence $\carrier[X]\subseteq\unions{\ran{D}}$
        by \obv{subseteq}.
    Follows by \obv{subseteq_antisymmetric}.
\end{proof}
